\documentclass{article}

\usepackage[final]{neurips_2026}

\usepackage[utf8]{inputenc} % allow utf-8 input
\usepackage[T1]{fontenc}    % use 8-bit T1 fonts
\usepackage{hyperref}       % hyperlinks
\usepackage{url}            % simple URL typesetting
\usepackage{booktabs}       % professional-quality tables
\usepackage{amsfonts}       % blackboard math symbols
\usepackage{nicefrac}       % compact symbols for 1/2, etc.
\usepackage{microtype}      % microtypography
\usepackage{xcolor}         % colors

\title {AI and Text Analytics:\\Mining Insights from Customer Feedback}


% The \author macro works with any number of authors. There are two commands
% used to separate the names and addresses of multiple authors: \And and \AND.
%
% Using \And between authors leaves it to LaTeX to determine where to break the
% lines. Using \AND forces a line break at that point. So, if LaTeX puts 3 of 4
% authors names on the first line, and the last on the second line, try using
% \AND instead of \And before the third author name.


\author{ MSc Data Science 2026\\ University of Bristol}
  

\begin{document}


\maketitle


\begin{abstract}
Abstract (0.25 pages): a brief summary of your project, stating the main challenge, the ideas you
tested, and key findings from the experiments.
\end{abstract} 

\section{Introduction}
Introduction (1.5 pages): outline the task you have been assigned, its motivations, the available
data and requirements for a good solution. You may wish to include examples from the dataset
and/or some data exploration plots.

\subsection{Background and motivation}
\subsection{Project outline}
\subsection{Initial data exploration}


\section{Methods}
Methods (3.5 pages):
a) Explain the baseline system, including important preprocessing steps and chosen text
representations or features. Include figures and/or equations to make your design clearer.
b) Discuss the axes of experimentation, how they could affect performance.
c) Describe your proposed improvements along each chosen axis. Include hypotheses about why you think these improvements will help.
\subsection{Preprocessing and pipeline overview}
\subsection{Axis 1: Text representation}
\subsection{Axis 2: Topic modeling}

\section{Experimental Evaluation}
Experimental evaluation (3.5 pages):
a) Explain your choice of performance metrics and their limitations.
b) Describe your experimental setup: how each dataset split was used, how hyperparameters
were set, and any other important information about how your methods were run. It is not
necessary to list all the software packages you used, but reference any libraries used for the
core parts of your system (e.g., links to HuggingFace models).
c) Present your results in suitable plots and/or tables, providing interpretation of the results to
highlight important observations about the comparative performance of your approaches.
d) Analyse the errors that your methods make to understand if there are particular text
patterns that cause more frequent errors.
e) Interpret and discuss your results, explaining the strengths and limitations of each method, and insights or recommendations based on your analysis.


\subsection{Experimental setup}
\subsection{Axis 1: Text Representation}
\subsubsection{Metrics and limitations}
\subsubsection{Results and discussion}


\subsection{Axis 2: Topic Modeling}
\subsubsection{Metrics and limitations}
\subsubsection{Results and discussion}
\subsection{Pipeline comparison}
(how can improvements be made based on actionable insights gatheres?)



\section{Conclusions}
\label{others}
Conclusions (1 page): describe the key findings from your experiments, explaining what worked
well, what did not work, and the open challenges for this task.
\subsection{Key findings}
\subsection{Challenges}


\section*{References}


References follow the acknowledgments in the camera-ready paper. Use unnumbered first-level heading for
the references. Any choice of citation style is acceptable as long as you are
consistent. It is permissible to reduce the font size to \verb+small+ (9 point)
when listing the references.
Note that the Reference section does not count towards the page limit.
\medskip


{
\small


[1] Alexander, J.A.\ \& Mozer, M.C.\ (1995) Template-based algorithms for
connectionist rule extraction. In G.\ Tesauro, D.S.\ Touretzky and T.K.\ Leen
(eds.), {\it Advances in Neural Information Processing Systems 7},
pp.\ 609--616. Cambridge, MA: MIT Press.


[2] Bower, J.M.\ \& Beeman, D.\ (1995) {\it The Book of GENESIS: Exploring
  Realistic Neural Models with the GEneral NEural SImulation System.}  New York:
TELOS/Springer--Verlag.


[3] Hasselmo, M.E., Schnell, E.\ \& Barkai, E.\ (1995) Dynamics of learning and
recall at excitatory recurrent synapses and cholinergic modulation in rat
hippocampal region CA3. {\it Journal of Neuroscience} {\bf 15}(7):5249-5262.
}

\end{document}